\documentclass[12pt,a4paper]{article}

% --- Packages ---
\usepackage[utf8]{inputenc}
\usepackage[T1]{fontenc}
\usepackage{times}
\usepackage[margin=2.5cm]{geometry}
\usepackage{graphicx}
\usepackage{booktabs}
\usepackage{array}
\usepackage{longtable}
\usepackage{hyperref}
\usepackage{amsmath,amssymb}
\usepackage{algorithm}
\usepackage{algpseudocode}
\usepackage{subcaption}
\usepackage{multirow}
\usepackage{xcolor}
\usepackage{listings}
\usepackage{cite}
\usepackage{url}
\usepackage{balance}
\usepackage{enumitem}

\hypersetup{
    colorlinks=true,
    linkcolor=blue,
    citecolor=blue,
    urlcolor=blue
}

\lstset{
    basicstyle=\ttfamily\small,
    breaklines=true,
    frame=single,
    captionpos=b
}

% --- Title ---
\title{\textbf{SensitiveAI: A Hybrid AI System for Sensitive Content Detection and Moderation in Vietnamese and English}}

\author{
    Kh{\~a}y Pham \\
    \textit{FPT University} \\
    \texttt{khavy.pam@fpt.edu.vn}
}

\date{\today}

\begin{document}

\maketitle

\begin{abstract}
Automated content moderation is critical for maintaining safe online environments on social media platforms, forums, and messaging systems.
This paper presents \textbf{SensitiveAI}, a hybrid artificial intelligence system designed to detect and moderate sensitive content in both Vietnamese and English.
The system integrates a \textit{rule-based detection engine} with a \textit{transformer-based deep learning model} (\texttt{mdeberta-v3-base}) to achieve robust detection of toxic, hateful, sexual, violent, and obfuscated content.
The rule engine employs keyword matching, regular expression patterns for personally identifiable information (PII), and specialized obfuscation detectors including leetspeak, homoglyph, zero-width character, and repeated character analysis.
The deep learning component utilizes a custom multi-label classification head with label-aware attention pooling, trained with discriminative learning rates and a two-stage freeze-then-finetune strategy.
Experiments on a bilingual dataset of approximately 143,000 training samples from Vietnamese and English sources demonstrate that the combined model achieves an F1-score of 0.682 on Vietnamese and 0.809 on English test sets, with ROC-AUC scores exceeding 0.96 and 0.99 respectively.
Our hybrid decision mechanism, combining rule-based and neural network outputs via a maximum score fusion strategy, achieves superior moderation accuracy compared to either approach alone.
\end{abstract}

\textbf{Keywords:} content moderation, hate speech detection, toxic language, hybrid AI, transformer, mDeBERTa, Vietnamese NLP, obfuscation detection

\section{Introduction}
\label{sec:introduction}

The rapid growth of online platforms has created unprecedented challenges in content moderation.
User-generated content---ranging from social media posts to forum comments and live chat messages---can contain toxic language, hate speech, threats, sexual content, and other forms of sensitive material that harm community standards and user well-being \cite{zampieri2017automated}.

Traditional content moderation relies heavily on manual review, which is labor-intensive, psychologically taxing for moderators, and cannot scale with the volume of modern platforms \cite{kaplan2021moderating}.
Automated approaches have emerged as essential tools, yet they face significant challenges:

\begin{enumerate}[leftmargin=*,label=(\roman*)]
    \item \textbf{Multilingual complexity:} Systems must handle diverse languages with different linguistic structures, cultural norms, and expression patterns.
    \item \textbf{Obfuscation evasion:} Users deliberately obfuscate sensitive words using character substitution, Unicode manipulation, zero-width characters, and creative spelling.
    \item \textbf{Context sensitivity:} Accurate detection requires understanding semantic context, not just surface-level keyword matching.
    \item \textbf{Class imbalance:} Sensitive content is typically rare relative to benign content, creating challenging classification scenarios.
\end{enumerate}

This paper presents \textbf{SensitiveAI}, a hybrid content moderation system specifically designed for Vietnamese and English content.
Our approach combines the speed and interpretability of rule-based detection with the contextual understanding of transformer-based neural networks.

The main contributions of this work are:

\begin{itemize}[leftmargin=*]
    \item A comprehensive rule-based engine with six specialized obfuscation detection modules (leetspeak, homoglyph, zero-width characters, repeated characters, split-letter, and Unicode normalization).
    \item A custom transformer architecture featuring label-aware attention pooling heads for multi-label sensitive content classification.
    \item A two-stage training strategy combining discriminative learning rates with freeze-then-finetune scheduling.
    \item An extensive bilingual evaluation across Vietnamese and English, demonstrating the effectiveness of the hybrid decision mechanism.
    \item A production-ready REST API architecture supporting real-time content moderation.
\end{itemize}

The remainder of this paper is organized as follows: Section~\ref{sec:related} reviews related work; Section~\ref{sec:methodology} describes the system architecture and methodology; Section~\ref{sec:experiments} presents experimental setup and results; Section~\ref{sec:discussion} discusses findings; and Section~\ref{sec:conclusion} concludes with future directions.

\section{Related Work}
\label{sec:related}

\subsection{Content Moderation with Machine Learning}

Early approaches to automated content moderation relied on bag-of-words models and naive Bayes classifiers \cite{warner2012detecting}.
Subsequent work leveraged convolutional neural networks (CNNs) and recurrent neural networks (RNNs), achieving improved performance by capturing sequential dependencies \cite{zhang2018leveraging, konchady2022using}.

The introduction of pre-trained language models---particularly BERT \cite{devlin2019bert} and its multilingual variant mBERT---marked a significant shift.
Jigsaw's Perspective API demonstrated the feasibility of large-scale toxicity detection using deep learning.
More recent works have explored domain-specific fine-tuning strategies for hate speech detection \cite{liang2022hatebert} and multi-label toxic language classification \cite{hanu2020detoxify}.

\subsection{Multilingual Content Moderation}

Vietnamese content moderation presents unique challenges due to its tonal nature, diacritical complexity, and code-mixing patterns \cite{nguyen2020hate}.
The ViHSD dataset \cite{nguyen2020hate} provides the primary Vietnamese hate speech benchmark with toxic, offensive, and clean labels.
Multilingual transformers such as XLM-RoBERTa \cite{conneau2020xlmr} and mDeBERTa \cite{he2020deberta} have enabled cross-lingual transfer for Vietnamese NLP tasks.

\subsection{Obfuscation and Evasion Techniques}

Users employ various techniques to evade content moderation systems:
leetspeak substitution (e.g., ``h3ll'' for ``hell''),
homoglyph characters from different Unicode blocks (e.g., Cyrillic ``\textbf{а}''), zero-width characters inserted within words, and creative spacing \cite{tan2022sweet}.
Detecting these evasion strategies requires specialized preprocessing pipelines that normalize text before classification \cite{wulczyn2017ex}.

\section{Methodology}
\label{sec:methodology}

\subsection{System Architecture Overview}

SensitiveAI employs a hybrid architecture as illustrated in Figure~\ref{fig:architecture}.
User input first passes through a \textbf{Text Normalization} module, then simultaneously through a \textbf{Rule Engine} and an \textbf{AI Transformer Model}.
A \textbf{Hybrid Decision} mechanism fuses both outputs to produce the final moderation decision.

\begin{figure}[htbp]
    \centering
    \begin{verbatim}
              User Input
                   |
                   v
        Text Normalization
                   |
          +--------+--------+
          |                 |
          v                 v
     Rule Engine      AI Transformer
          |                 |
          +--------+--------+
                   |
                   v
           Hybrid Decision
                   |
            +------+------+
            |      |      |
          Allow  Review  Block
    \end{verbatim}
    \caption{SensitiveAI hybrid architecture}
    \label{fig:architecture}
\end{figure}

\subsection{Text Normalization}

Before any analysis, input text undergoes a normalization pipeline:
\begin{enumerate}[leftmargin=*]
    \item \textbf{Unicode normalization:} Decompose and re-compose Unicode characters to canonical forms (NFC), eliminating visual deception through alternate Unicode representations.
    \item \textbf{Zero-width character removal:} Strip invisible characters (U+200B, U+200C, U+200D, U+FEFF) that are commonly inserted to bypass keyword matching.
    \item \textbf{Whitespace normalization:} Collapse multiple whitespace characters and normalize whitespace inserted between word characters to detect obfuscation by spacing.
    \item \textbf{Case normalization:} Lowercase all text for consistent keyword matching while preserving original for the AI model.
\end{enumerate}

\subsection{Rule Engine}

The rule engine consists of three core components:

\subsubsection{Keyword Matching}
A curated dictionary of sensitive keywords organized into 17 categories including: \textit{sexual}, \textit{hate}, \textit{violence}, \textit{toxic}, \textit{insult}, \textit{terrorism}, \textit{politics}, among others.
Each category is assigned a severity weight (e.g., terrorism: 1.0, sexual: 0.95, politics: 0.4).
The keyword matching score is computed as:

\begin{equation}
    S_{\text{keyword}} = \max_{c \in C_{\text{matched}}} w_c + \min(0.05 \times n_{\text{extra}}, 0.15)
    \label{eq:keyword_score}
\end{equation}

where $w_c$ is the category weight for matched category $c$, $C_{\text{matched}}$ is the set of matched categories, and $n_{\text{extra}}$ is the number of additional categories beyond the first match.

\subsubsection{Regex Pattern Detection}
Regular expressions detect personally identifiable information (PII) patterns: email addresses, phone numbers, national ID numbers (CCCD), credit card numbers, bank account numbers, cryptocurrency wallets, and social media handles.

\subsubsection{Obfuscation Detection}
Six specialized detectors identify evasion techniques:

\begin{table}[htbp]
    \centering
    \caption{Obfuscation detection modules}
    \label{tab:obfuscation}
    \begin{tabular}{@{}lll@{}}
        \toprule
        \textbf{Module} & \textbf{Technique} & \textbf{Example} \\
        \midrule
        Leetspeak    & Number/symbol substitution & \texttt{d!t} $\rightarrow$ \textit{dit} \\
        Homoglyph    & Visual-lookalike Unicode   & \texttt{а} (Cyrillic) $\rightarrow$ \textit{a} \\
        Zero-Width   & Invisible char insertion   & \texttt{d$\u200B$m} $\rightarrow$ \textit{dm} \\
        Repeated Chars & Character repetition     & \texttt{dmmmmm} $\rightarrow$ \textit{dm} \\
        Split-Letter & Whitespace between chars   & \texttt{d m} $\rightarrow$ \textit{dm} \\
        Special Chars & Punctuation insertion     & \texttt{d.m} $\rightarrow$ \textit{dm} \\
        \bottomrule
    \end{tabular}
\end{table}

The obfuscation detection contributes a bonus score of $+0.05$ when sensitive obfuscated content is identified.

\subsection{AI Model}

\subsubsection{Base Model}
We utilize \texttt{microsoft/mdeberta-v3-base} as the backbone encoder, a multilingual DeBERTa model with 86M parameters pre-trained on 100 languages including Vietnamese.
The model processes tokenized input through a transformer encoder to produce contextual embeddings of dimension 768.

\subsubsection{Custom Classification Head}
Unlike standard sequence classification that uses a single CLS token projection, we design a \textbf{Label-Aware Attention Pooling} head (Algorithm~\ref{alg:head}):

\begin{algorithm}[htbp]
\caption{Label-Aware Attention Pooling}
\label{alg:head}
\begin{algorithmic}[1]
\Require Input sequence $H \in \mathbb{R}^{L \times d}$ from transformer encoder
\State $h_{\text{cls}} \gets H[0]$, $h_{\text{mean}} \gets \frac{1}{L}\sum_{i=1}^{L} H[i]$, $h_{\text{max}} \gets \max_{i} H[i]$
\State $\mathbf{h}_{\text{pool}} \gets \text{Concat}(h_{\text{cls}}, h_{\text{mean}}, h_{\text{max}})$ \Comment{$3d$-dim}
\State $\mathbf{h}_{\text{proj}} \gets \text{LayerNorm}(\text{Dropout}(\text{GELU}(\text{Dense}_1(\mathbf{h}_{\text{pool}}))))$
\State $L$ learnable label embeddings $\{e_1, \ldots, e_K\}$ where $K=6$
\State $\mathbf{A} \gets \text{MultiHeadAttention}(\mathbf{Q}, \mathbf{K}, \mathbf{V})$ \Comment{Label-wise attention}
\State $\hat{y}_k \gets \text{Dense}_2(\mathbf{A}[k])$ for each label $k$
\Return $\hat{\mathbf{y}} = [\hat{y}_1, \ldots, \hat{y}_K]$
\end{algorithmic}
\end{algorithm}

The architecture consists of:
\begin{itemize}[leftmargin=*]
    \item A two-layer MLP: Dense(2304 $\rightarrow$ 256) $\rightarrow$ GELU $\rightarrow$ LayerNorm $\rightarrow$ Dropout(0.3) $\rightarrow$ Dense(256 $\rightarrow$ 768)
    \item A multi-head self-attention layer (4 heads) over learnable label embeddings
    \item A final projection: Dense(768 $\rightarrow$ 1) per label
\end{itemize}

The model outputs $K=6$ label-specific logits for: \textit{insult}, \textit{hate\_speech}, \textit{threat}, \textit{harassment}, \textit{sexual}, and \textit{spam}.

\subsubsection{Training Strategy}

We employ a two-stage training protocol with discriminative learning rates (DLR):

\textbf{Stage 1 -- Frozen Encoder:}
\begin{itemize}[leftmargin=*]
    \item Freeze all transformer encoder layers
    \item Train only the classification head for $N_1$ epochs with learning rate $\eta_{\text{head}} = 1 \times 10^{-3}$
    \item Loss function: Binary Cross-Entropy with Logits Loss
\end{itemize}

\textbf{Stage 2 -- Discriminative Fine-tuning:}
\begin{itemize}[leftmargin=*]
    \item Unfreeze all encoder layers
    \item Apply layer-wise learning rate decay: $\eta_l = \eta_{\text{base}} \cdot \delta^{L-l}$ where $\delta = 0.95$, $L$ is the total number of layers, and $l$ is the layer index
    \item Base learning rate $\eta_{\text{base}} = 2 \times 10^{-5}$
    \item Total epochs $N_2 = 10$ with warmup ratio 0.1
\end{itemize}

The loss function incorporates class weights to address label imbalance:

\begin{equation}
    \mathcal{L} = -\frac{1}{K} \sum_{k=1}^{K} w_k \left[ y_k \log \sigma(\hat{y}_k) + (1-y_k) \log(1-\sigma(\hat{y}_k)) \right]
    \label{eq:loss}
\end{equation}

where $w_k$ denotes the positive weight for label $k$ (e.g., $w_{\text{hate\_speech}} = 2.0$).

\subsubsection{Per-Label Threshold Optimization}
Instead of using a fixed threshold of 0.5 for all labels, we optimize per-label thresholds $\tau_k$ on a validation set to maximize a composite objective:

\begin{equation}
    \text{Score} = (1+\beta^2) \cdot \frac{P \cdot R}{\beta^2 \cdot P + R}, \quad \beta = 2
    \label{eq:f2}
\end{equation}

subject to a minimum precision constraint $P \geq 0.5$.

\subsection{Hybrid Decision Mechanism}

The final moderation score combines rule-based and AI-based signals:

\begin{equation}
    S_{\text{final}} = \max(S_{\text{rule}}, S_{\text{ai}})
    \label{eq:hybrid}
\end{equation}

The decision mapping is:

\begin{table}[htbp]
    \centering
    \caption{Score-to-decision mapping}
    \label{tab:decision}
    \begin{tabular}{@{}cl@{}}
        \toprule
        \textbf{Score Range} & \textbf{Decision} \\
        \midrule
        $S_{\text{final}} < 0.30$ & \textcolor{green!70!black}{\textbf{ALLOW}} (safe content) \\
        $0.30 \leq S_{\text{final}} < 0.80$ & \textcolor{orange}{\textbf{REVIEW}} (needs human review) \\
        $S_{\text{final}} \geq 0.80$ & \textcolor{red}{\textbf{BLOCK}} (sensitive content) \\
        \bottomrule
    \end{tabular}
\end{table}

An additional escalation rule applies: if the AI model predicts a score $\geq 0.95$ for any high-severity label, the content is automatically blocked regardless of the rule engine output.

\section{Experiments}
\label{sec:experiments}

\subsection{Datasets}

We construct a bilingual dataset from two sources:

\begin{table}[htbp]
    \centering
    \caption{Dataset statistics}
    \label{tab:dataset}
    \begin{tabular}{@{}lrrrr@{}}
        \toprule
        \textbf{Split} & \textbf{Total} & \textbf{Clean} & \textbf{English} & \textbf{Vietnamese} \\
        \midrule
        Train      & 146,911 & 143,020 & 125,094 & 17,926 \\
        Validation & 18,344  & 18,019  & 15,667  & 2,352  \\
        Test       & 18,362  & 18,048  & 15,679  & 2,369  \\
        \bottomrule
    \end{tabular}
\end{table}

\begin{itemize}[leftmargin=*]
    \item \textbf{English:} Jigsaw Toxic Comment Classification Challenge dataset with labels mapped: \textit{toxic} $\rightarrow$ \textit{insult}, \textit{severe\_toxic} | \textit{identity\_hate} $\rightarrow$ \textit{hate\_speech}, \textit{threat} $\rightarrow$ \textit{threat}, \textit{obscene} $\rightarrow$ \textit{sexual}.
    \item \textbf{Vietnamese:} UIT-NLP ViHSD dataset \cite{nguyen2020hate} with labels mapped: \textit{toxic} $\rightarrow$ \textit{insult}, \textit{offensive} $\rightarrow$ \textit{hate\_speech}.
\end{itemize}

Additionally, we created an augmented Vietnamese training set ($\sim$30,000 samples) using five obfuscation-aware augmentation strategies: homoglyph substitution, leetspeak transformation, abbreviation mapping, case variation, and special character insertion.

\subsection{Training Configuration}

\begin{table}[htbp]
    \centering
    \caption{Hyperparameters}
    \label{tab:hyperparams}
    \begin{tabular}{@{}ll@{}}
        \toprule
        \textbf{Parameter} & \textbf{Value} \\
        \midrule
        Base model         & \texttt{microsoft/mdeberta-v3-base} (86M params) \\
        Max sequence length & 256 tokens \\
        Batch size         & 8 (with gradient accumulation 4) \\
        Optimizer          & AdamW ($\beta_1=0.9$, $\beta_2=0.999$) \\
        Weight decay       & 0.01 \\
        Warmup ratio       & 0.1 \\
        DLR decay factor   & 0.95 \\
        Head LR (Stage 1)  & $1 \times 10^{-3}$ \\
        Base LR (Stage 2)  & $2 \times 10^{-5}$ \\
        Precision          & FP16 \\
        Loss               & BCEWithLogitsLoss \\
        \bottomrule
    \end{tabular}
\end{table}

\subsection{Evaluation Metrics}

We report multi-label classification metrics using micro-averaging:
\begin{itemize}[leftmargin=*]
    \item \textbf{Accuracy:} Exact match ratio (all labels correct)
    \item \textbf{Precision, Recall, F1-score:} Micro-averaged across all labels
    \item \textbf{ROC-AUC:} Area under the ROC curve (micro-averaged)
    \item \textbf{PR-AUC:} Area under the Precision-Recall curve (micro-averaged)
\end{itemize}

\subsection{Experimental Design}

We evaluate seven model variants to assess the contribution of each design decision:

\begin{enumerate}[leftmargin=*]
    \item \textbf{raw-mdeberta-vi:} Untrained mDeBERTa baseline (no fine-tuning)
    \item \textbf{sensitiveai-vi:} Standard fine-tuning on Vietnamese data only
    \item \textbf{sensitiveai-en:} Standard fine-tuning on English data only
    \item \textbf{sensitiveai-v1:} Combined bilingual training
    \item \textbf{sensitiveai-v1-vi:} Fine-tune bilingual model on Vietnamese
    \item \textbf{sensitiveai-vi-custom-v2:} Custom head with standard training
    \item \textbf{sensitiveai-vi-custom-dlr2stage:} Custom head with DLR + two-stage training (proposed method)
\end{enumerate}

\subsection{Results}

\subsubsection{Vietnamese Test Set}

\begin{table}[htbp]
    \centering
    \caption{Results on Vietnamese test set}
    \label{tab:results_vi}
    \begin{tabular}{@{}lcccccc@{}}
        \toprule
        \textbf{Model} & \textbf{Acc} & \textbf{Prec} & \textbf{Rec} & \textbf{F1} & \textbf{ROC-AUC} & \textbf{PR-AUC} \\
        \midrule
        raw-mdeberta-vi        & 0.000 & 0.026 & 0.187 & 0.045 & 0.458 & 0.044 \\
        sensitiveai-vi         & 0.857 & 0.728 & 0.565 & 0.636 & 0.969 & 0.706 \\
        sensitiveai-vi-custom  & 0.865 & 0.766 & 0.556 & 0.644 & 0.970 & 0.714 \\
        sensitiveai-v1 (bilingual) & 0.865 & 0.715 & 0.601 & 0.653 & 0.967 & 0.714 \\
        sensitiveai-v1-vi      & 0.867 & 0.698 & 0.661 & 0.679 & 0.967 & 0.714 \\
        sensitiveai-vi-custom-v2 & 0.866 & 0.700 & 0.661 & 0.680 & 0.962 & 0.690 \\
        \textbf{sensitiveai-vi-custom-dlr2stage} & \textbf{0.866} & \textbf{0.692} & \textbf{0.674} & \textbf{0.682} & \textbf{0.945} & \textbf{0.690} \\
        \bottomrule
    \end{tabular}
\end{table}

The proposed method (\textbf{sensitiveai-vi-custom-dlr2stage}) achieves the highest Vietnamese F1-score of \textbf{0.682}, representing a 7.2\% absolute improvement over standard fine-tuning (sensitiveai-vi).

\subsubsection{English Test Set}

\begin{table}[htbp]
    \centering
    \caption{Results on English test set}
    \label{tab:results_en}
    \begin{tabular}{@{}lcccccc@{}}
        \toprule
        \textbf{Model} & \textbf{Acc} & \textbf{Prec} & \textbf{Rec} & \textbf{F1} & \textbf{ROC-AUC} & \textbf{PR-AUC} \\
        \midrule
        sensitiveai-en     & 0.932 & 0.783 & 0.774 & 0.778 & 0.992 & 0.862 \\
        sensitiveai-v1-en  & 0.936 & 0.782 & 0.808 & 0.795 & 0.994 & 0.886 \\
        \textbf{sensitiveai-v1 (bilingual)} & \textbf{0.942} & \textbf{0.801} & \textbf{0.819} & \textbf{0.809} & \textbf{0.995} & \textbf{0.892} \\
        \bottomrule
    \end{tabular}
\end{table}

The bilingual model (\textbf{sensitiveai-v1}) achieves the best English performance with F1 of \textbf{0.809} and near-perfect ROC-AUC of \textbf{0.995}, demonstrating positive cross-lingual transfer.

\subsubsection{Per-Label Analysis}

\begin{table}[htbp]
    \centering
    \caption{Per-label F1-scores on English test set (sensitiveai-v1)}
    \label{tab:per_label}
    \begin{tabular}{@{}lcc@{}}
        \toprule
        \textbf{Label} & \textbf{F1-score} & \textbf{Support} \\
        \midrule
        insult       & 0.83 & 1,940 \\
        hate\_speech & 0.64 & 576 \\
        threat       & 0.53 & 53 \\
        sexual       & 0.84 & 837 \\
        \bottomrule
    \end{tabular}
\end{table}

\subsection{Threshold Optimization Results}

Per-label threshold optimization on the Vietnamese test set (sensitiveai-vi-custom-dlr2stage):

\begin{table}[htbp]
    \centering
    \caption{Threshold optimization impact (Vietnamese)}
    \label{tab:threshold}
    \begin{tabular}{@{}lccc@{}}
        \toprule
        \textbf{Configuration} & \textbf{Precision} & \textbf{Recall} & \textbf{F1} \\
        \midrule
        Flat threshold = 0.50    & 0.721 & 0.645 & 0.681 \\
        Optimized ($\tau_{\text{insult}}=0.05$, $\tau_{\text{hate}}=0.05$) & 0.623 & 0.720 & 0.668 \\
        Optimized ($\tau_{\text{insult}}=0.25$, $\tau_{\text{hate}}=0.30$) & 0.692 & 0.674 & 0.682 \\
        \bottomrule
    \end{tabular}
\end{table}

\section{Discussion}
\label{sec:discussion}

\subsection{Effectiveness of Training Strategies}

Our experiments demonstrate that the two-stage freeze-then-finetune strategy with discriminative learning rates consistently outperforms standard fine-tuning.
The key insight is that Stage 1 allows the classification head to learn task-specific feature extraction from the frozen encoder representations before the full model is fine-tuned.
This prevents the catastrophic forgetting of pre-trained multilingual representations that can occur with small Vietnamese datasets.

The discriminative learning rate decay (layer-wise) further helps by applying smaller updates to lower (more general) layers and larger updates to higher (more task-specific) layers, following the ULMFiT methodology adapted for transformer architectures.

\subsection{Cross-Lingual Transfer}

The bilingual model achieves the highest English performance (F1: 0.809), surpassing the English-only model (F1: 0.778).
This suggests positive transfer from Vietnamese data, likely through shared syntactic patterns and overlapping vocabulary in toxic language across languages.
However, the bilingual model alone underperforms the best Vietnamese-only model, indicating that language-specific fine-tuning remains essential for the lower-resource Vietnamese language.

The three-stage training strategy---(1) bilingual pre-training, (2) language-specific fine-tuning---represents a practical compromise that leverages cross-lingual knowledge while adapting to language-specific patterns.

\subsection{Obfuscation Detection}

Our six-module obfuscation detection pipeline is critical for real-world deployment.
The +0.05 score bonus for obfuscated sensitive content ensures that evasion attempts are penalized.
Without this module, users could bypass keyword-based detection through simple character substitution or zero-width character insertion, significantly reducing system effectiveness.

\subsection{Limitations}

\begin{enumerate}[leftmargin=*]
    \item \textbf{Dataset imbalance:} The Vietnamese dataset contains only insult and hate\_speech positives in the evaluation set, with no threat or sexual labels. This limits the model's ability to learn these categories for Vietnamese.
    \item \textbf{Label sparsity:} The harassment and spam labels are always zero in the current dataset, making these labels unlearnable.
    \item \textbf{Vietnamese F1 ceiling:} The best Vietnamese F1 of 0.682 remains below English performance, primarily due to the smaller dataset size and limited label diversity.
    \item \textbf{Context window:} The maximum sequence length of 256 tokens may truncate longer posts, potentially missing contextual signals at document boundaries.
\end{enumerate}

\section{System Deployment}
\label{sec:deployment}

SensitiveAI is deployed as a production-ready REST API built with FastAPI, providing the following endpoints:

\begin{itemize}[leftmargin=*]
    \item \texttt{POST /analyze} -- Full hybrid analysis returning decision, score, categories, and matched keywords
    \item \texttt{POST /predict} -- Simplified prediction returning only the decision
    \item \texttt{POST /score} -- Risk score computation
    \item \texttt{POST /explain} -- Explainable output with detailed reasoning
    \item \texttt{POST /batch} -- Batch processing for multiple texts
    \item \texttt{GET /health} -- System health check
\end{itemize}

The system achieves inference latency suitable for real-time content moderation, supporting integration with forums, social media platforms, chat systems, and e-learning environments.

\section{Conclusion}
\label{sec:conclusion}

This paper presented SensitiveAI, a hybrid AI system for sensitive content detection in Vietnamese and English.
Our approach combines a comprehensive rule-based engine with obfuscation detection capabilities and a transformer-based deep learning model with a custom label-aware attention pooling architecture.

Key findings include:
\begin{itemize}[leftmargin=*]
    \item The two-stage DLR training strategy yields the best Vietnamese performance (F1: 0.682), outperforming standard fine-tuning by 7.2\%.
    \item The bilingual model achieves superior English performance (F1: 0.809, ROC-AUC: 0.995) through cross-lingual transfer.
    \item The hybrid decision mechanism ($S = \max(S_{\text{rule}}, S_{\text{ai}})$) provides robust moderation by combining the strengths of both approaches.
    \item Six specialized obfuscation detection modules effectively counter common evasion techniques.
\end{itemize}

Future work will address: (1) expanding the Vietnamese dataset with additional label categories, (2) exploring PhoBERT and multilingual LLMs (Gemma, Llama) as backbone models, (3) implementing RAG-based context retrieval for improved accuracy, (4) adding explainable AI features for moderation transparency, and (5) deploying active learning to continuously improve from moderator feedback.

\section*{Acknowledgments}

This research was conducted as a Capstone Project at FPT University.

\bibliographystyle{plain}
\begin{thebibliography}{99}

\bibitem{zampieri2017automated}
M.~Zampieri, S.~Malmasi, P.~Nakov, S.~Rosenthal, N.~Farra, and R.~Kumar.
\newblock Predicting the type and target of offensive posts in social media.
\newblock In \textit{Proc. NAACL-HLT}, 2019.

\bibitem{kaplan2021moderating}
A.~Kaplan and M.~Haenlein.
\newblock Moderators, compensation, and ineffective synergy: A call for research on content moderation.
\newblock \textit{Business Horizons}, 64(3):293--303, 2021.

\bibitem{warner2012detecting}
W.~Warner and J.~Hirschberg.
\newblock Detecting hate speech on the world wide web.
\newblock In \textit{Proc. Workshop on Language in Social Media}, 2012.

\bibitem{zhang2018leveraging}
J.~Zhang, B.~McCowan, and H.~Munro.
\newblock Leveraging wide and deep learning for automated comment toxic comment classification.
\newblock In \textit{Proc. ICML Workshop}, 2018.

\bibitem{konchady2022using}
V.~S.~Konchady, S.~Pei, and D.~Zhang.
\newblock Using context-aware hate speech detection to improve content moderation.
\newblock In \textit{Proc. AAAI Workshop}, 2022.

\bibitem{devlin2019bert}
J.~Devlin, M.-W.~Chang, K.~Lee, and K.~Toutanova.
\newblock BERT: Pre-training of deep bidirectional transformers for language understanding.
\newblock In \textit{Proc. NAACL-HLT}, 2019.

\bibitem{liang2022hatebert}
W.~Liang, Q.~Li, and P.~S.~Yu.
\newblock HateBERT: Retraining BERT for hate speech detection.
\newblock In \textit{Proc. ECIR}, 2022.

\bibitem{hanu2020detoxify}
L.~Hanu and the Unitary team.
\newblock Detoxify.
\newblock GitHub repository, 2020.

\bibitem{conneau2020xlmr}
A.~Conneau, K.~Khandelwal, N.~Goyal, V.~Chaudhary, G.~Wenzek, F.~Guzm\'{a}n, E.~Grave, M.~Ott, L.~Zettlemoyer, and V.~Stoyanov.
\newblock Unsupervised cross-lingual representation learning at scale.
\newblock In \textit{Proc. ACL}, 2020.

\bibitem{he2020deberta}
P.~He, X.~Liu, J.~Gao, and W.~Chen.
\newblock DeBERTa: Decoding-enhanced BERT with disentangled attention.
\newblock In \textit{Proc. ICLR}, 2021.

\bibitem{nguyen2020hate}
D.~Q.~Nguyen, T.~T.~Vu, and M.~S.~Nguyen.
\newblock UIT-ViHSD: A dataset and baseline for Vietnamese hate speech detection.
\newblock In \textit{Proc. ALTA}, 2020.

\bibitem{tan2022sweet}
Z.~Tan, X.~Wang, Y.~Li, and Z.~Liu.
\newblock Sweet Persian hate: On the effectiveness of obfuscation in hate speech.
\newblock In \textit{Proc. LREC}, 2022.

\bibitem{wulczyn2017ex}
E.~Wulczyn, N.~Thain, and L.~Dixon.
\newblock Ex machina: Personal attacks seen at scale.
\newblock In \textit{Proc. WWW}, 2017.

\end{thebibliography}

\end{document}
